\textbf{Phần 1. Trắc nghiệm lựa chọn}
\setcounterpageref{ex}{0}

\begin{ex} %[4D5V26]
Tính chất giao hoán của phép cộng là
\choice
{$a - b = b - a$}
{$(a + b) + c = a + (b + c)$}
{\True $a + b = b + a$}
{$a + 0 = a$}
\loigiai{
$a + b = b + a$.\\
\textbf{Đáp án của câu hỏi là C}
}
\end{ex}

\begin{ex} %[4D5V26]
Kết quả của phép tính $450\,000 + 350\,000$ là
\choice
{$700\,000$}
{\True $800\,000$}
{$900\,000$}
{$850\,000$}
\loigiai{
$450\,000 + 350\,000 = 800\,000$.\\
\textbf{Đáp án của câu hỏi là B}
}
\end{ex}

\begin{ex} %[4D5V26]
Hiệu của $1\,000\,000$ và $250\,000$ là
\choice
{$850\,000$}
{$650\,000$}
{$700\,000$}
{\True $750\,000$}
\loigiai{
$1\,000\,000 - 250\,000 = 750\,000$.\\
\textbf{Đáp án của câu hỏi là D}
}
\end{ex}

\begin{ex} %[4D5V26]
Tìm hai số biết tổng là $80$ và hiệu là $14$. Số lớn là
\choice
{\True $47$}
{$44$}
{$50$}
{$52$}
\loigiai{
$(80 + 14) : 2 = 47$.\\
\textbf{Đáp án của câu hỏi là A}
}
\end{ex}

\begin{ex} %[4D5V26]
Tìm $x$ biết: $x + 45\,000 = 120\,000$.
\choice
{$65\,000$}
{$85\,000$}
{\True $75\,000$}
{$70\,000$}
\loigiai{
$x = 120\,000 - 45\,000 = 75\,000$.\\
\textbf{Đáp án của câu hỏi là C}
}
\end{ex}

\begin{ex} %[4D5V26]
Tính thuận tiện: $250 + 184 + 750$.
\choice
{$1084$}
{\True $1184$}
{$1284$}
{$1194$}
\loigiai{
$(250 + 750) + 184 = 1000 + 184 = 1184$.\\
\textbf{Đáp án của câu hỏi là B}
}
\end{ex}

\begin{ex} %[4D5V26]
Hai thùng dầu chứa tất cả $240\text{ l}$. Thùng thứ nhất chứa nhiều hơn thùng thứ hai $40\text{ l}$. Số lít dầu ở thùng thứ hai là
\choice
{$90\text{ l}$}
{$110\text{ l}$}
{$120\text{ l}$}
{\True $100\text{ l}$}
\loigiai{
$(240 - 40) : 2 = 100\text{ l}$.\\
\textbf{Đáp án của câu hỏi là D}
}
\end{ex}

\begin{ex} %[4D5V26]
Một khu đất hình chữ nhật có chu vi $160\text{ m}$, chiều dài hơn chiều rộng $20\text{ m}$. Diện tích khu đất đó là
\choice
{\True $1500\text{ m}^2$}
{$1200\text{ m}^2$}
{$1600\text{ m}^2$}
{$1800\text{ m}^2$}
\loigiai{
Nửa chu vi: $80\text{ m}$.\\
Chiều dài: $(80 + 20) : 2 = 50\text{ m}$.\\
Chiều rộng: $80 - 50 = 30\text{ m}$.\\
Diện tích: $50 \times 30 = 1500\text{ m}^2$.\\
\textbf{Đáp án của câu hỏi là A}
}
\end{ex}

\begin{ex} %[4D5V26]
Tính nhanh giá trị biểu thức: $12\,500 + 34\,800 + 87\,500 + 65\,200$.
\choice
{$190\,000$}
{$210\,000$}
{\True $200\,000$}
{$220\,000$}
\loigiai{
$(12\,500 + 87\,500) + (34\,800 + 65\,200) = 100\,000 + 100\,000 = 200\,000$.\\
\textbf{Đáp án của câu hỏi là C}
}
\end{ex}

\begin{ex} %[4D5V26]
Một nhà máy đợt 1 may được $45\,000$ chiếc áo, đợt 2 may được $52\,000$ chiếc áo. Để đạt chỉ tiêu $120\,000$ chiếc áo thì đợt 3 nhà máy cần may số chiếc áo là
\choice
{$22\,000$ chiếc}
{\True $23\,000$ chiếc}
{$24\,000$ chiếc}
{$25\,000$ chiếc}
\loigiai{
$120\,000 - (45\,000 + 52\,000) = 23\,000$ chiếc.\\
\textbf{Đáp án của câu hỏi là B}
}
\end{ex}

\begin{ex} %[4D5V26]
Tuổi của mẹ và Lan cộng lại là $42$ tuổi. Mẹ sinh Lan khi mẹ $26$ tuổi. Tuổi của mẹ hiện nay là
\choice
{$32$ tuổi}
{$36$ tuổi}
{$38$ tuổi}
{\True $34$ tuổi}
\loigiai{
Mẹ hơn Lan 26 tuổi.\\
Tuổi mẹ: $(42 + 26) : 2 = 34$ tuổi.\\
\textbf{Đáp án của câu hỏi là D}
}
\end{ex}

\begin{ex} %[4D5V26]
Tìm hai số chẵn có tổng bằng $210$, biết giữa chúng có đúng $4$ số chẵn khác. Số lớn là
\choice
{\True $110$}
{$110$}
{$112$}
{$114$}
\loigiai{
Giữa chúng có 4 số chẵn khác nên hiệu hai số là $(4 + 1) \times 2 = 10$.\\
Số lớn là: $(210 + 10) : 2 = 110$.\\
\textbf{Đáp án của câu hỏi là A}
}
\end{ex}

\textbf{Phần 2. Trắc nghiệm đúng sai}
\setcounterpageref{ex}{0}

\begin{ex} %[4D5V26]
Xét tính đúng sai của các phép toán:
\choiceTF
{\True $125\,000 + 875\,000 = 1\,000\,000$.}
{$2\,000\,000 - 850\,000 = 1\,250\,000$.}
{\True $2\,000\,000 - 850\,000 = 1\,150\,000$.}
{\True $(a + b) + c = a + (b + c)$.}
\loigiai{
\begin{itemize}
    \item $125\,000 + 875\,000 = 1\,000\,000$. Đúng.
    \item $2\,000\,000 - 850\,000 = 1\,150\,000 \neq 1\,250\,000$. Ý 2 Sai, Ý 3 Đúng.
    \item Tính chất kết hợp: Đúng.
\end{itemize}
\begin{itemize}
    \item a) \textbf{Đúng}.
    \item b) \textbf{Sai}.
    \item c) \textbf{Đúng}.
    \item d) \textbf{Đúng}.
\end{itemize}
}
\end{ex}

\begin{ex} %[4D5V26]
Hai bao gạo có tất cả $120\text{ kg}$. Bao 1 nặng hơn bao 2 là $20\text{ kg}$. Xét tính đúng sai:
\choiceTF
{\True Bao 1 nặng $70\text{ kg}$.}
{\True Bao 2 nặng $50\text{ kg}$.}
{Nếu lấy ở bao 1 ra $10\text{ kg}$ thì hai bao bằng nhau.}
{\True Nếu lấy ở bao 1 chuyển sang bao 2 là $10\text{ kg}$ thì hai bao bằng nhau.}
\loigiai{
\begin{itemize}
    \item Bao 1: $(120 + 20) : 2 = 70\text{ kg}$. Đúng.
    \item Bao 2: $120 - 70 = 50\text{ kg}$. Đúng.
    \item Lấy ở bao 1 bỏ đi 10 kg thì bao 1 còn 60 kg, bao 2 là 50 kg (vẫn không bằng nhau). Sai.
    \item Chuyển 10 kg sang bao 2 thì cả 2 bao đều có 60 kg. Đúng.
\end{itemize}
\begin{itemize}
    \item a) \textbf{Đúng}.
    \item b) \textbf{Đúng}.
    \item c) \textbf{Sai}.
    \item d) \textbf{Đúng}.
\end{itemize}
}
\end{ex}

\textbf{Phần 3. Trả lời ngắn}
\setcounterpageref{ex}{0}

\begin{ex} %[4D5V26]
Tính tổng: $1\,250\,000 + 3\,750\,000$.
\shortans{5000000}
\loigiai{
\textbf{Đáp số:} $5000000$\\
$5\,000\,000$.
}
\end{ex}

\begin{ex} %[4D5V26]
Hai số có tổng là $96$, hiệu là $24$. Tìm số lớn.
\shortans{60}
\loigiai{
\textbf{Đáp số:} $60$\\
$(96 + 24) : 2 = 60$.
}
\end{ex}

\begin{ex} %[4D5V26]
Tính thuận tiện: $450 + 890 + 550$.
\shortans{1890}
\loigiai{
\textbf{Đáp số:} $1890$\\
$(450 + 550) + 890 = 1890$.
}
\end{ex}

\begin{ex} %[4D5V26]
Hai số lẻ có tổng là $180$, biết giữa chúng có $2$ số lẻ khác. Tìm số bé.
\shortans{87}
\loigiai{
\textbf{Đáp số:} $87$\\
Hiệu: $(2 + 1) \times 2 = 6 \Rightarrow (180 - 6) : 2 = 87$.
}
\end{ex}

\begin{ex} %[4D5V26]
Một hình chữ nhật có nửa chu vi là $60\text{ cm}$, chiều dài hơn chiều rộng $16\text{ cm}$. Tính diện tích hình chữ nhật theo xăng-ti-mét vuông.
\shortans{836}
\loigiai{
\textbf{Đáp số:} $836$\\
Dài: $38\text{ cm}$, Rộng: $22\text{ cm} \Rightarrow 38 \times 22 = 836\text{ cm}^2$.
}
\end{ex}

\begin{ex} %[4D5V26]
Tính tổng: $100 + 200 + 300 + \dots + 1000$.
\shortans{5500}
\loigiai{
\textbf{Đáp số:} $5500$\\
$(100 + 1000) \times 10 : 2 = 5500$.
}
\end{ex}
